\documentclass[12pt]{article}
\usepackage[T1]{fontenc}
\usepackage[utf8]{inputenc}
\usepackage{lmodern}
\usepackage{textcomp}
\usepackage{enumitem}
\usepackage{geometry}
\geometry{margin=1in}
\begin{document}
\begin{center}
Answer Key - Political Science - Undergraduate Level Assessment 24 \
\small (Answers are ordered exactly as in the question paper.)
\end{center}

\section*{Multiple Choice}
\begin{enumerate}[label=\arabic*.]
\item \textbf{Answer:} A
\\ \textit{Options:} A) contribute money to candidates for election; B) coordinate local get-out-the-vote campaigns; C) promote the defeat of incumbents in the federal and state legislatures; D) organize protest demonstrations and other acts of civil disobedience
\\ \textit{Note:} Correct option: A - contribute money to candidates for election
\item \textbf{Answer:} D
\\ \textit{Options:} A) Ways and Means; B) Judiciary; C) Ethics; D) Rules
\\ \textit{Note:} Correct option: D - Rules
\item \textbf{Answer:} A
\\ \textit{Options:} A) A federal mandate; B) A constitutional amendment; C) Affirmative action; D) Tort reform
\\ \textit{Note:} Correct option: A - A federal mandate
\item \textbf{Answer:} B
\\ \textit{Options:} A) It encouraged greater involvement in European politics; B) It was designed as the antithesis of European politics; C) It created a large standing army; D) It encouraged the centralization of political power in the US
\\ \textit{Note:} Correct option: B - It was designed as the antithesis of European politics
\item \textbf{Answer:} A
\\ \textit{Options:} A) Withdrawal of economic trade rights with the domestic market.; B) Export controls protecting technological advantage and further foreign policy objectives.; C) Control of munitions and arms sales.; D) Import restrictions to protect a domestic market from foreign goods.
\\ \textit{Note:} Correct option: A - Withdrawal of economic trade rights with the domestic market.
\end{enumerate}

\section*{True / False}
\begin{enumerate}[label=\arabic*.]
\setcounter{enumi}{5}
\item \textbf{Answer:} False
\\ \textit{Options:} A) True; B) False
\\ \textit{Note:} LLM-generated false statement. Correct answer: 'The importance of the issue to a particular individual or group'.
\item \textbf{Answer:} True
\\ \textit{Options:} A) True; B) False
\\ \textit{Note:} LLM-generated true statement based on MCQ.
\item \textbf{Answer:} False
\\ \textit{Options:} A) True; B) False
\\ \textit{Note:} LLM-generated false statement. Correct answer: 'Autocracy'.
\item \textbf{Answer:} False
\\ \textit{Options:} A) True; B) False
\\ \textit{Note:} LLM-generated false statement. Correct answer: 'Voluntary reliance on an external power for security'.
\item \textbf{Answer:} False
\\ \textit{Options:} A) True; B) False
\\ \textit{Note:} LLM-generated false statement. Correct answer: 'Voters have become more focused on individual candidates.'.
\end{enumerate}

\section*{Long Form Response}
\begin{enumerate}[label=\arabic*.]
\setcounter{enumi}{10}
\item \textbf{Answer:} Embrace multi-polarity, show greater restraint internationally and require other states to meet their own security burdens.. Explain why this is the correct answer and provide supporting evidence. Reference key facts or concepts that support this choice.
\\ \textit{Note:} Long-form derived from MCQ; award credit for stating the correct idea and giving supporting reasoning.
\item \textbf{Answer:} less clearly identified with consistent political ideologies. Explain why this is the correct answer and provide supporting evidence. Reference key facts or concepts that support this choice.
\\ \textit{Note:} Long-form derived from MCQ; award credit for stating the correct idea and giving supporting reasoning.
\end{enumerate}

\vfill
\noindent\textit{End of Answer Key}
\end{document}